\HMTNarrativeHeading{From determinantal structure to orientation}

The Planck constants are presented after the objects that determine their
action sections. The dodecaphasic register contains the local Hadamard block;
its integral quotient has a calculable order. The determinantal reader
composes this datum with the coordinates already generated and fixes a
decimal valuation. The scale and coefficient of a section perform different
functions, which the proof keeps explicit.

The return distinguishes the sections before and after the turn. Their
coefficients preserve their origin in the character and the regional
continuation. The expression of action in physical units is introduced
through the corresponding dimensional chart. The construction thus
transmits a section, its coefficient, its scale, and the reading rule.
The subsequent use of the constant preserves these relations.

Angle and counter-angle are developed next as orientation coordinates.
The first reading shifts the blocks of alpha in the completion base; the
second composes the dimensionless return coefficient with alpha. The
circular chart expresses the result in degrees or radians and preserves
the number of turns when the lift is used. This sequence explains why
the angular coordinates are important intermediaries between action and
the subsequent geometrical readers.

The conjugate channels receive this orientation and, together with marked
incidence, evaluate the Barbero--Immirzi functional. The electronic
realization uses its trajectory and its return readers, and receives the
quotient of the action sections. The narrative preserves the order of its
antecedents and distinguishes the dimensionless coefficient from the
magnitude of action it helps to represent. The same constant can act at
several stages without losing the map through which it was obtained.

This development prepares a concrete physical composition. Action and clock
determine an energy; thermal transduction relates this energy to
information; the area reader incorporates marked incidence. The
longitudinal composition retains the inverse, its memory, and the
inscribed incidences, determining a length on the declared radial
reference. Its arrival modulus relates energy, conjugate length, and
reduced radius; the two normalized identities fix the common coupling.
The article preserves the calculations of these compositions to show
how the arithmetic and geometrical results participate in a common physical
reading.
